% !TEX root = ../../ctfp-print.tex

\lettrine[lhang=0.17]{今}{日、関数型プログラミングを語るには}モナドについて言及せずにはいられない。しかし、Eugenio Moggi がモナドではなく Lawvere 理論に注目していたという別の宇宙が存在する。その宇宙を探索してみよう。

\section{普遍代数}

代数を様々な抽象レベルで記述する方法は多くある。モノイド、群、環といったものを記述するための汎用言語を見つけようとしている。最も単純なレベルでは、これらの構成はすべて集合の要素に対する\emph{演算}と、それらの演算が満たすべき\emph{法則}を定義する。例えば、モノイドは結合的な二項演算によって定義できる。また単位元と単位法則もある。しかし少し想像力を使えば、単位元を零項演算—引数を取らず集合の特定の要素を返す演算—に変換できる。群について語りたければ、元を取りその逆元を返す単項演算子を加える。それに対応する左・右逆元法則がある。環は二つの二項演算子といくつかの法則を定義する。以下同様だ。

大きな絵としては、代数は様々な値の $n$ に対する $n$ 項演算の集合と、等式恒等式の集合によって定義される。これらの恒等式はすべて全称量化されている。結合性の等式は三つの元のすべての可能な組み合わせについて満たされなければならない。以下同様だ。

ちなみに、これにより体（field）は考察から除外される。単純な理由として、（加法に関する）単位元であるゼロが乗法に関する逆元を持たないからだ。体に対する逆元法則は全称量化できない。

普遍代数のこの定義は、演算（関数）を射に置き換えることで $\Set$ 以外の圏に拡張できる。集合の代わりに、対象 $a$（汎用対象と呼ばれる）を選択する。単項演算は $a$ の自己射に過ぎない。しかし、他のアリティについてはどうだろうか（\newterm{アリティ}（arity）は与えられた演算の引数の数だ）？二項演算（アリティ2）は積 $a\times{}a$ から $a$ への射として定義できる。一般的な $n$ 項演算は $a$ の $n$ 乗から $a$ への射だ：
\[\alpha_n \Colon a^n \to a\]
零項演算は終対象（$a$ の零乗）からの射だ。したがって、任意の代数を定義するために必要なのは、一つの特殊な対象 $a$ の冪を対象とする圏だけだ。特定の代数はこの圏のホム集合にエンコードされる。これが Lawvere 理論の要点だ。

Lawvere 理論の導出は多くのステップを経るので、ここにロードマップを示す：

\begin{enumerate}
  \tightlist
  \item
        有限集合の圏 $\cat{FinSet}$。
  \item
        そのスケルトン $\cat{F}$。
  \item
        その反対 $\Fop$。
  \item
        Lawvere 理論 $\cat{L}$：圏 $\cat{Law}$ における対象。
  \item
        Lawvere 圏のモデル $M$：圏\\
        $\cat{Mod}(\cat{L}, \Set)$ における対象。
\end{enumerate}

\begin{figure}[H]
  \centering
  \includegraphics[width=0.8\textwidth]{images/lawvere1.png}
\end{figure}

\section{Lawvere 理論}

すべての Lawvere 理論は共通の骨格を共有する。Lawvere 理論のすべての対象は、積（実際には冪）を使って一つの対象から生成される。しかし、一般的な圏でこれらの積をどのように定義するのか？積はより単純な圏からの写像を使って定義できることがわかる。実際、このより単純な圏は積の代わりに余積を定義してもよく、\emph{反変}関手を使ってそれらを目標の圏に埋め込む。反変関手は余積を積に変換し、単射を射影に変換する。

Lawvere 圏の骨格として自然な選択は、有限集合の圏 $\cat{FinSet}$ だ。この圏には空集合 $\varnothing$、一元集合 $1$、二元集合 $2$ などが含まれる。この圏のすべての対象は余積を使って一元集合から生成できる（空集合を零項余積の特殊な場合として扱う）。例えば、二元集合は二つの一元集合の和であり、$2 = 1 + 1$ と表される。Haskell では：

\src{snippet01}
しかし、空集合は一つしかないと考えるのが自然でも、多くの異なる一元集合が存在しうる。特に、集合 $1 + \varnothing$ は集合 $\varnothing + 1$ とは異なり、また $1$ とも異なる—それらすべてが同型であっても。集合の圏における余積は結合的でない。すべての同型な集合を同一視する圏を構築することでこの状況を改善できる。そのような圏は\newterm{スケルトン}（skeleton）と呼ばれる。言い換えれば、すべての Lawvere 理論の骨格は $\cat{FinSet}$ のスケルトン $\cat{F}$ だ。この圏の対象は自然数（ゼロを含む）と同一視でき、これは $\cat{FinSet}$ の元の数に対応する。余積は加法の役割を果たす。$\cat{F}$ の射は有限集合間の関数に対応する。例えば、$\varnothing$ から $n$ への一意な射がある（空集合が始対象であるため）、$n$ から $\varnothing$ への射はない（$\varnothing \to \varnothing$ を除いて）、$1$ から $n$ への射は $n$ 個（単射）、$n$ から $1$ への射は一つ、などだ。ここで $n$ は $\cat{F}$ における対象を表し、$\cat{FinSet}$ の同型を通じて同一視されたすべての $n$ 元集合に対応する。

圏 $\cat{F}$ を使って、\newterm{Lawvere 理論}（Lawvere theory）を特殊な関手を装備した圏 $\cat{L}$ として形式的に定義できる：
\[I_{\cat{L}} \Colon \Fop \to \cat{L}\]
この関手は対象上の全単射でなければならず、有限積を保存しなければならない（$\Fop$ における積は $\cat{F}$ における余積と同じだ）：
\[I_{\cat{L}} (m\times{}n) = I_{\cat{L}} m\times{}I_{\cat{L}} n\]
この関手は対象上の恒等写像として特徴づけられることがある。つまり、$\cat{F}$ と $\cat{L}$ の対象は同じだということだ。そのため、それらに同じ名前を使う—自然数で表す。ただし、$\cat{F}$ の対象は集合と同じではないことを念頭に置いておくこと（それらは同型な集合のクラスだ）。

$\cat{L}$ のホム集合は、一般に $\Fop$ のものより豊かだ。これらは $\cat{FinSet}$ の関数に対応するもの以外の射を含むことがある（後者は時に\newterm{基本積演算}（basic product operations）と呼ばれる）。Lawvere 理論の等式法則はそれらの射にエンコードされている。

重要な観察は、$\cat{F}$ における一元集合 $1$ が $\cat{L}$ においても $1$ と呼ぶ対象に写され、$\cat{L}$ の他のすべての対象が自動的にこの対象の冪になることだ。例えば、$\cat{F}$ における二元集合 $2$ は余積 $1 + 1$ であるため、$\cat{L}$ においては積 $1 \times 1$（または $1^2$）に写されなければならない。この意味で、圏 $\cat{F}$ は $\cat{L}$ の対数のように振る舞う。

$\cat{L}$ の射の中には、関手 $I_{\cat{L}}$ によって $\cat{F}$ から転送されたものがある。それらは $\cat{L}$ において構造的な役割を果たす。特に余積単射 $i_k$ は積射影 $p_k$ になる。有益な直感は、射影：
\[p_k \Colon 1^n \to 1\]
を $n$ 個の変数の関数の雛形として想像することだ。第 $k$ 番目の変数を除くすべてを無視する。逆に、$\cat{F}$ における定数射 $n \to 1$ は $\cat{L}$ における対角射 $1 \to 1^n$ になる。これらは変数の複製に対応する。

$\cat{L}$ において興味深い射は、射影以外の $n$ 項演算を定義するものだ。一つの Lawvere 理論を別のものと区別するのはそれらの射だ。代数を定義する乗算、加算、単位元の選択などがそれだ。しかし、$\cat{L}$ を完全な圏にするために、複合演算 $n \to m$（あるいは同等に $1^n \to 1^m$）も必要だ。圏の単純な構造のため、それらは $n \to 1$ 型のより単純な射の積になることがわかる。これは積を返す関数は関数の積であるという命題の一般化だ；以前見たように、ホム関手は連続だ。

\begin{figure}[H]
  \centering
  \includegraphics[width=0.8\textwidth]{images/lawvere1.png}
  \caption{Lawvere 理論 $\cat{L}$ は $\Fop$ に基づいており、積を定義する「退屈な」射を継承する。$n$ 項演算を記述する「興味深い」射（点線矢印）を追加する。}
\end{figure}

Lawvere 理論は圏 $\cat{Law}$ を形成し、そこでの射は有限積を保存し、関手 $I$ と可換な関手だ。二つのそのような理論 $(\cat{L}, I_{\cat{L}})$ と $(\cat{L'}, I'_{\cat{L'}})$ が与えられると、それらの間の射は次の条件を満たす関手 $F \Colon \cat{L} \to \cat{L'}$ だ：
\begin{gather*}
  F (m \times n) = F m \times F n \\
  F \circ I_{\cat{L}} = I'_{\cat{L'}}
\end{gather*}
Lawvere 理論間の射は、一つの理論を別の理論の中で解釈するというアイデアをカプセル化する。例えば、逆元を無視すれば、群の乗算をモノイドの乗算として解釈できる。

Lawvere 圏の最も単純な自明な例は $\Fop$ 自身だ（$I_{\cat{L}}$ として恒等関手を選択することに対応する）。演算も法則も持たないこの Lawvere 理論は $\cat{Law}$ における始対象だ。

この時点で、非自明な Lawvere 理論の例を示すことが非常に役立つだろうが、最初にモデルが何であるかを理解しなければ説明が難しい。

\section{Lawvere 理論のモデル}

Lawvere 理論を理解する鍵は、一つのそのような理論が同じ構造を共有する多くの個別の代数を一般化することを認識することだ。例えば、モノイドの Lawvere 理論はモノイドであることの本質を記述する。それはすべてのモノイドに対して有効でなければならない。特定のモノイドはそのような理論のモデルになる。モデルは Lawvere 理論 $\cat{L}$ から集合の圏 $\Set$ への関手として定義される。（モデルに他の圏を使う Lawvere 理論の一般化もあるが、ここでは $\Set$ だけに集中する。）$\cat{L}$ の構造は積に大きく依存しているので、そのような関手が有限積を保存することを要求する。$\cat{L}$ のモデル—Lawvere 理論 $\cat{L}$ 上の代数とも呼ばれる—は、次のような関手によって定義される：
\begin{gather*}
  M \Colon \cat{L} \to \Set \\
  M (a \times b) \cong M a \times M b
\end{gather*}
積の保存は\emph{同型を除いて}のみ要求することに注意する。これは非常に重要で、積の厳密な保存はほとんどの興味深い理論を除外してしまう。

モデルによる積の保存は、$\Set$ における $M$ の像が集合 $M 1$—$\cat{L}$ からの対象 $1$ の像—の冪によって生成される集合の列であることを意味する。この集合を $a$ と呼ぼう。（この集合は時に\emph{ソート}（sort）と呼ばれ、そのような代数は\newterm{単ソート}（single-sorted）と呼ばれる。Lawvere 理論を多ソート代数に拡張する一般化も存在する。）特に、$\cat{L}$ からの二項演算は関数に写される：
\[a \times a \to a\]
任意の関手と同様に、$\cat{L}$ の複数の射が $\Set$ の同じ関数に潰される可能性がある。

ちなみに、すべての法則が全称量化された等式であるという事実は、すべての Lawvere 理論が自明なモデルを持つことを意味する：すべての対象を一元集合に写し、すべての射をそれに対する恒等関数に写す定数関手だ。

$\cat{L}$ における $m \to n$ の形の一般的な射は、次の関数に写される：
\[a^m \to a^n\]
二つの異なるモデル $M$ と $N$ があれば、それらの間の自然変換は $n$ によってインデックスされた関数の族だ：
\[\mu_n \Colon M n \to N n\]
あるいは同等に：
\[\mu_n \Colon a^n \to b^n\]
ここで $b = N 1$ だ。

自然性条件が $n$ 項演算の保存を保証することに注意する：
\[N f \circ \mu_n = \mu_1 \circ M f\]
ここで $f \Colon n \to 1$ は $\cat{L}$ における $n$ 項演算だ。

モデルを定義する関手はモデルの圏 $\cat{Mod}(\cat{L}, \Set)$ を形成し、射として自然変換を持つ。

自明な Lawvere 圏 $\Fop$ のモデルを考えてみよう。そのようなモデルは $1$ における値 $M 1$ によって完全に決定される。$M 1$ は任意の集合にできるので、$\Set$ の集合と同じ数だけこれらのモデルが存在する。さらに、$\cat{Mod}(\Fop, \Set)$ における各射（関手 $M$ と $N$ の間の自然変換）は $1$ における成分によって一意に決定される。逆に、すべての関数 $M 1 \to N 1$ は二つのモデル $M$ と $N$ の間の自然変換を誘導する。したがって $\cat{Mod}(\Fop, \Set)$ は $\Set$ と同値だ。

\section{モノイドの理論}

Lawvere 理論の最も単純な非自明な例は、モノイドの構造を記述するものだ。これは可能なすべてのモノイドの構造を蒸留した単一の理論であり、この理論のモデルがモノイドの圏 $\cat{Mon}$ 全体にわたる。以前、すべてのモノイドは射の適切なサブセットを同一視することで適切な自由モノイドから得られることを示す\hyperref[free-monoids]{普遍構成}を見た。したがって、単一の自由モノイドがすでに多くのモノイドを一般化する。しかし、自由モノイドは無限に多い。モノイドの Lawvere 理論 $\cat{L}_{\cat{Mon}}$ はそれらすべてを一つのエレガントな構成にまとめる。

すべてのモノイドは単位元を持たなければならないので、$\cat{L}_{\cat{Mon}}$ に $0$ から $1$ へ進む特殊な射 $\eta$ が必要だ。$\cat{F}$ には対応する射が存在できないことに注意する。そのような射は逆方向、つまり $1$ から $0$ へ向かうことになり、$\cat{FinSet}$ では一元集合から空集合への関数となる。そのような関数は存在しない。

次に、$\cat{L}_{\cat{Mon}}(2, 1)$ の元である射 $2 \to 1$ を考えよう。これはすべての二項演算の雛形を含まなければならない。$\cat{Mod}(\cat{L}_{\cat{Mon}}, \Set)$ においてモデルを構成するとき、これらの射はデカルト積 $M 1 \times M 1$ から $M 1$ への関数に写される。言い換えれば、二引数の関数だ。

問いはこうだ：モノイド演算だけを使って、どれだけの二引数の関数を実装できるか。二つの引数を $a$ と $b$ と呼ぼう。両方の引数を無視してモノイド単位元を返す関数が一つある。次に、それぞれ $a$ と $b$ を返す二つの射影がある。続いて $ab$、$ba$、$aa$、$bb$、$aab$ などを返す関数がある……実際、生成元 $a$ と $b$ を持つ自由モノイドの元の数と同じだけの二引数の関数がある。$\cat{L}_{\cat{Mon}}(2, 1)$ はそれらすべての射を含まなければならないことに注意する。なぜなら、モデルの一つが自由モノイドであるからだ。自由モノイドではそれらは異なる関数に対応する。他のモデルでは $\cat{L}_{\cat{Mon}}(2, 1)$ の複数の射が一つの関数に潰される可能性があるが、自由モノイドではそうならない。

$n$ 個の生成元を持つ自由モノイドを $n^*$ と表すと、ホム集合 $\cat{L}(2, 1)$ をモノイドの圏 $\cat{Mon}$ におけるホム集合 $\cat{Mon}(1^*, 2^*)$ と同一視できる。一般に、$\cat{L}_{\cat{Mon}}(m, n)$ を $\cat{Mon}(n^*, m^*)$ とする。言い換えれば、圏 $\cat{L}_{\cat{Mon}}$ は自由モノイドの圏の反対圏だ。

モノイドの Lawvere 理論のモデルの圏 $\cat{Mod}(\cat{L}_{\cat{Mon}}, \Set)$ はすべてのモノイドの圏 $\cat{Mon}$ と同値だ。

\section{Lawvere 理論とモナド}

覚えているように、代数的理論はモナドを使って記述できる—特に\hyperref[algebras-for-monads]{モナドの代数}を使って。そのため、Lawvere 理論とモナドの間に繋がりがあっても不思議ではない。

まず、Lawvere 理論がどのようにモナドを誘導するかを見てみよう。それは忘却関手と自由関手の間の\hyperref[free-forgetful-adjunctions]{随伴}を通じて行われる。忘却関手 $U$ は各モデルに集合を割り当てる。この集合は $\cat{Mod}(\cat{L}, \Set)$ からの関手 $M$ を $\cat{L}$ の対象 $1$ で評価することで得られる。

$U$ を導出するもう一つの方法は、$\Fop$ が $\cat{Law}$ における始対象であるという事実を利用することだ。これは、任意の Lawvere 理論 $\cat{L}$ に対して、$\Fop \to \cat{L}$ という一意な関手が存在することを意味する。この関手はモデル上の反対の関手を誘導する（モデルは理論\emph{から}集合への関手であるため）：
\[\cat{Mod}(\cat{L}, \Set) \to \cat{Mod}(\Fop, \Set)\]
しかし、議論したように、$\Fop$ のモデルの圏は $\Set$ と同値なので、忘却関手が得られる：
\[U \Colon \cat{Mod}(\cat{L}, \Set) \to \Set\]
このように定義された $U$ は常に左随伴、すなわち自由関手 $F$ を持つことが示せる。

これは有限集合に対して容易にわかる。自由関手 $F$ は自由代数を生成する。自由代数は有限個の生成元 $n$ から生成された $\cat{Mod}(\cat{L}, \Set)$ における特定のモデルだ。$F$ を表現可能関手として実装できる：
\[\cat{L}(n, -) \Colon \cat{L} \to \Set\]
これが確かに自由であることを示すには、忘却関手への左随伴であることを示すだけでよい：
\[\cat{Mod}(\cat{L}, \Set)(\cat{L}(n, -), M) \cong \Set(n, U(M))\]
右辺を単純化しよう：
\[\Set(n, U(M)) \cong \Set(n, M 1) \cong (M 1)^n \cong M n\]
（射の集合が指数と同型であるという事実を使った。この場合、それは単なる反復積だ。）随伴は米田の補題の結果だ：
\[[\cat{L}, \Set](\cat{L}(n, -), M) \cong M n\]
忘却関手と自由関手が合わさって、$\Set$ 上の\hyperref[monads-categorically]{モナド} $T = U \circ F$ を定義する。したがって、すべての Lawvere 理論はモナドを生成する。

\hyperref[algebras-for-monads]{このモナドの代数}の圏はモデルの圏と同値であることがわかる。

モナド代数は、モナドを使って形成された式を評価する方法を定義することを思い出してほしい。Lawvere 理論は式を生成するために使える $n$ 項演算を定義する。モデルはこれらの式を評価する手段を提供する。

しかし、モナドと Lawvere 理論の繋がりは双方向ではない。有限項モナドだけが Lawvere 理論に繋がる。\newterm{有限項モナド}（finitary monad）は有限項関手に基づく。$\Set$ 上の有限項関手は、有限集合上の作用によって完全に決定される。任意の集合 $a$ 上の作用は次のコエンドを使って評価できる：
\[F a = \int^n a^n \times (F n)\]
コエンドが余積または和を一般化するので、この公式は冪級数展開の一般化だ。あるいは、関手は一般化されたコンテナであるという直感を使えばよい。その場合、$a$ の有限項コンテナは形の和と内容として記述できる。ここで $F n$ は $n$ 個の要素を格納するための形の集合であり、内容は $n$ タプルの要素、それ自体 $a^n$ の元だ。例えば、リスト（関手として）は有限項であり、各アリティに対して一つの形を持つ。ツリーはアリティごとにより多くの形を持つ、などだ。

まず、Lawvere 理論から生成されたすべてのモナドは有限項であり、コエンドとして表現できる：
\[T_{\cat{L}} a = \int^n a^n \times \cat{L}(n, 1)\]
逆に、$\Set$ 上の任意の有限項モナド $T$ が与えられると、Lawvere 理論を構成できる。$T$ のクライスリ圏を構成することから始める。覚えているように、クライスリ圏における $a$ から $b$ への射は、基底圏における射によって与えられる：
\[a \to T b\]
有限集合に限定すると：
\[m \to T n\]
このクライスリ圏の反対圏 $\cat{Kl}^\mathit{op}_{T}$ は、有限集合に限定すると、問題の Lawvere 理論だ。特に、$\cat{L}$ における $n$ 項演算を記述するホム集合 $\cat{L}(n, 1)$ は、ホム集合 $\cat{Kl}_{T}(1, n)$ によって与えられる。

プログラミングで遭遇するほとんどのモナドは有限項であることがわかる。継続モナドが注目すべき例外だ。有限項演算を超えて Lawvere 理論の概念を拡張することも可能だ。

\section{コエンドとしてのモナド}

コエンドの公式をより詳しく探索しよう。
\[T_{\cat{L}} a = \int^n a^n \times \cat{L}(n, 1)\]
まず、このコエンドは $\cat{F}$ における次のように定義されたプロ関手 $P$ にわたって取られる：
\[P n m = a^n \times \cat{L}(m, 1)\]
このプロ関手は第一引数 $n$ において反変だ。射をどのようにリフトするかを考えよう。$\cat{FinSet}$ における射は有限集合の写像 $f \Colon m \to n$ だ。そのような写像は $n$ 元集合から $m$ 個の要素の選択を記述する（重複も許される）。これは $a$ の冪の写像にリフトできる。つまり（方向に注意）：
\[a^n \to a^m\]
リフトは単に $n$ 個の要素のタプル $(a_1, a_2, \ldots{}, a_n)$ から $m$ 個の要素を選択する（重複も許される）。

\begin{figure}[H]
  \centering
  \includegraphics[width=0.5\textwidth]{images/liftpower.png}
\end{figure}

\noindent
例えば、$f_k \Colon 1 \to n$—$n$ 元集合から第 $k$ 番目の要素を選択するもの—を取ってみよう。これは $a$ の $n$ タプルを取り、第 $k$ 番目を返す関数にリフトされる。

あるいは $f \Colon m \to 1$—すべての $m$ 個の要素を一つに写す定数関数—を取ってみよう。そのリフトは $a$ の単一の元を取り、それを $m$ 回複製する関数だ：
\[\lambda{}x \to (\underbrace{x, x, \ldots{}, x}_{m})\]
問題のプロ関手が第二引数において共変であることが直ちには明らかでないことに気づくかもしれない。ホム関手 $\cat{L}(m, 1)$ は実際には $m$ において反変だ。しかし、コエンドは圏 $\cat{L}$ ではなく圏 $\cat{F}$ において取る。コエンド変数 $n$ は有限集合（またはそのスケルトン）にわたる。圏 $\cat{L}$ は $\cat{F}$ の反対を含むので、$\cat{F}$ における射 $m \to n$ は $\cat{L}$ における $\cat{L}(n, m)$ の元だ（埋め込みは関手 $I_{\cat{L}}$ によって与えられる）。

$\cat{L}(m, 1)$ の $\cat{F}$ から $\Set$ への関手としての関手性を確認しよう。関数 $f \Colon m \to n$ をリフトしたい。したがって目標は $\cat{L}(m, 1)$ から $\cat{L}(n, 1)$ への関数を実装することだ。関数 $f$ に対応して、$\cat{L}$ には $n$ から $m$ への射がある（方向に注意）。この射を $\cat{L}(m, 1)$ と前合成すると $\cat{L}(n, 1)$ のサブセットが得られる。

\begin{figure}[H]
  \centering
  \begin{tikzcd}[column sep=large]
    \cat{L}(m, 1) \arrow[r] & \cat{L}(n, 1)\\
    {}^m \bullet \arrow[r, "f"'] & \bullet^n
  \end{tikzcd}
\end{figure}

\noindent
関数 $1 \to n$ をリフトすることで、$\cat{L}(1, 1)$ から $\cat{L}(n, 1)$ へ行けることに注意する。この事実は後で使う。

反変関手 $a^n$ と共変関手 $\cat{L}(m, 1)$ の積はプロ関手 $\Fop \times \cat{F} \to \Set$ だ。コエンドは、一部の元が同一視されるプロ関手のすべての対角成分の余積（非交和）として定義できることを思い出してほしい。同一視はコウェッジ条件に対応する。

ここで、コエンドはすべての $n$ にわたる集合 $a^n \times \cat{L}(n, 1)$ の非交和として始まる。同一視は\hyperref[ends-and-coends]{コエンドを余等化子として}表現することで生成できる。非対角項 $a^n \times \cat{L}(m, 1)$ から始める。対角に達するために、積の第一または第二成分に射 $f \Colon m \to n$ を適用できる。二つの結果は同一視される。

\begin{figure}[H]
  \centering
  \begin{tikzcd}
    & a^n \times \cat{L}(m, 1)
    \arrow[dl, "\langle f {,} \id \rangle"']
    \arrow[dr, "\langle \id {,} f \rangle"]
    & \\
    a^m \times \cat{L}(m, 1)
    & \scalebox{2.5}[1]{\sim}
    & a^n \times \cat{L}(n, 1) \\
    & f \Colon m \to n &
  \end{tikzcd}
\end{figure}

\noindent
先ほど、$f \Colon 1 \to n$ のリフトがこれら二つの変換をもたらすことを示した：
\[a^n \to a\]
と：
\[\cat{L}(1, 1) \to \cat{L}(n, 1)\]
したがって、$a^n \times \cat{L}(1, 1)$ から始めると、$\langle f, \id \rangle$ をリフトすると：
\[a \times \cat{L}(1, 1)\]
に達し、$\langle \id, f \rangle$ をリフトすると：
\[a^n \times \cat{L}(n, 1)\]
に達する。しかしこれは、$a^n \times \cat{L}(n, 1)$ のすべての元が $a \times \cat{L}(1, 1)$ と同一視できることを意味しない。なぜなら、$\cat{L}(n, 1)$ のすべての元が $\cat{L}(1, 1)$ から到達できるわけではないからだ。$\cat{F}$ からの射だけがリフトできることを思い出してほしい。$\cat{L}$ における非自明な $n$ 項演算は、射 $f \Colon 1 \to n$ をリフトすることで構成できない。

言い換えれば、$\cat{L}(n, 1)$ が基本射の適用を通じて $\cat{L}(1, 1)$ から到達できるコエンド公式内のすべての加数だけを同一視できる。それらはすべて $a \times \cat{L}(1, 1)$ と同値だ。基本射は $\cat{F}$ における射の像のものだ。

最も単純な Lawvere 理論の場合、$\Fop$ 自身において、これがどのように機能するかを見てみよう。そのような理論では、すべての $\cat{L}(n, 1)$ が $\cat{L}(1, 1)$ から到達できる。なぜなら $\cat{L}(1, 1)$ は恒等射だけを含む一元集合であり、$\cat{L}(n, 1)$ は $\cat{F}$ における単射 $1 \to n$ に対応する射のみを含んでいて、これらは\emph{基本射}だからだ。したがって、余積のすべての加数は同値であり、次が得られる：
\[T a = a \times \cat{L}(1, 1) = a\]
これは恒等モナドだ。

\section{副作用の Lawvere 理論}

モナドと Lawvere 理論の間にはこれほど強い繋がりがあるので、Lawvere 理論がプログラミングでモナドの代替として使えるかどうかという問いが自然に生じる。モナドの主な問題は、うまく合成できないことだ。モナド変換子を構築する汎用レシピは存在しない。Lawvere 理論はこの点で優位性がある：余積とテンソル積を使って合成できる。一方、有限項モナドだけが Lawvere 理論に簡単に変換できる。ここでの例外は継続モナドだ。この分野では現在も研究が続いている（参考文献を参照）。

Lawvere 理論が副作用を記述するためにどのように使えるかを示すために、伝統的に \code{Maybe} モナドを使って実装される例外の単純な場合を議論する。

\code{Maybe} モナドは、単一の零項演算 $0 \to 1$ を持つ Lawvere 理論によって生成される。この理論のモデルは $1$ をある集合 $a$ に写し、零項演算を次の関数に写す関手だ：

\src{snippet02}
コエンド公式を使って \code{Maybe} モナドを回復できる。零項演算の追加がホム集合 $\cat{L}(n, 1)$ に何をするかを考えてみよう。$\cat{L}(0, 1)$ に新しい元を作る（これは $\Fop$ には存在しない）だけでなく、$\cat{L}(n, 1)$ に新しい射も追加する。これらは型 $n \to 0$ の射を $0 \to 1$ と合成した結果だ。そのような寄与はすべてコエンド公式において $a^0 \times \cat{L}(0, 1)$ と同一視される。なぜなら次から得られるからだ：
\[a^n \times \cat{L}(0, 1)\]
$0 \to n$ を二つの異なる方法でリフトすることで。

\begin{figure}[H]
  \centering
  \begin{tikzcd}
    & a^n \times \cat{L}(0, 1)
    \arrow[dl, "\langle f {,} \id \rangle"']
    \arrow[dr, "\langle \id {,} f \rangle"]
    & \\
    a^0 \times \cat{L}(0, 1)
    & \scalebox{2.5}[1]{\sim}
    & a^n \times \cat{L}(n, 1) \\
    & f \Colon 0 \to n &
  \end{tikzcd}
\end{figure}

\noindent
コエンドは次に簡約される：
\[T_{\cat{L}} a = a^0 + a^1\]
あるいは、Haskell 記法を使うと：

\src{snippet03}
これは次と同値だ：

\src{snippet04}
この Lawvere 理論は例外の送出のみをサポートし、処理はサポートしないことに注意する。

\section{練習問題}

\begin{enumerate}
  \tightlist
  \item
        $\cat{F}$（$\cat{FinSet}$ のスケルトン）において $2$ と $3$ の間のすべての射を列挙せよ。
  \item
        モノイドの Lawvere 理論のモデルの圏がリストモナドのモナド代数の圏と同値であることを示せ。
  \item
        モノイドの Lawvere 理論はリストモナドを生成する。その二項演算が対応するクライスリ矢を使って生成できることを示せ。
  \item
        $\cat{FinSet}$ は $\Set$ の部分圏であり、それを $\Set$ に埋め込む関手がある。$\Set$ 上の任意の関手は $\cat{FinSet}$ に制限できる。有限項関手は自身の制限の左カン拡張であることを示せ。
\end{enumerate}

\section{参考文献}
\begin{enumerate}
  \tightlist
  \item
        \urlref{http://www.tac.mta.ca/tac/reprints/articles/5/tr5.pdf}{Functorial Semantics of Algebraic Theories}, F. William Lawvere
  \item
        \urlref{http://homepages.inf.ed.ac.uk/gdp/publications/Comp_Eff_Monads.pdf}{Notions of computation determine monads}, Gordon Plotkin and John Power
\end{enumerate}
